\section{Related Work}
\label{sec:related}

\paragraph{Acoustic monitoring for conservation.} Passive acoustic monitoring is
established for biodiversity and increasingly for threat detection (chainsaws,
gunshots, vehicles), with deployments such as Rainforest Connection converting
recordings into ranger alerts \citep{stowell2022bioacoustics}. These systems
predominantly use closed-set supervised classifiers and report in-domain metrics; cross-site generalization
under a fixed false-positive budget is rarely measured. A global expert assessment of
PAM practice identifies manual validation of detector output as a priority
bottleneck \citep{malerba2026pam}, which is why we treat reviewer attention,
not accuracy, as the operative constraint. Our negative
result quantifies the cross-site gap this reflects, and our detector targets the
label-scarce, open-taxonomy regime deployments actually face.


\paragraph{Few-shot bioacoustic event detection.} The DCASE few-shot bioacoustic
task formalizes detecting novel sound events from a handful of labeled examples
\citep{nolasco2023dcase}. Our prototype stage is in this spirit, since a class is defined by a few
verified positives, but we couple it to a background-only anomaly gate and an
explicit unknown option, so the system flags before any labels exist and
abstains when no prototype matches.


\paragraph{Open-set recognition and OOD detection.} Open-set recognition asks a
classifier to reject inputs outside its training classes; a large OOD-detection
literature scores novelty, including Mahalanobis distance in a deep feature
space \citep{lee2018mahalanobis}. We adopt the Mahalanobis-in-embedding scoring but invert the usual
framing: the ``in-distribution'' density is \emph{normal forest background}, not
the labeled classes, so the detector is trained without any positives and known
classes enter only as post-hoc prototypes.


\paragraph{Prototypical and metric learning.} Class-mean prototypes with cosine
similarity are standard in few-shot metric learning \citep{snell2017prototypical}. Our contribution is not the
prototype but its role in a deployable pipeline: prototypes are optional,
incrementally updatable from human confirmations, and gated behind a
background-calibrated anomaly threshold with reserved unknown mass.


\paragraph{Transfer of pretrained audio embeddings.} Large audio taggers pretrained on AudioSet \citep{gemmeke2017audioset}, such as
CNN14 \citep{kong2020panns}, yield general-purpose embeddings that transfer
across audio tasks, as does the bioacoustics-specific BirdNET
\citep{kahl2021birdnet}. We use one as a frozen $\phi$ and show that the choice of embedding,
not the detector head, governs cross-site deployability for forest threats.


\paragraph{Distribution shift, shortcut learning, and evaluation optimism.}
Models often exploit spurious, environment-correlated features rather than the
signal of interest \citep{geirhos2020shortcut}, and random train/test splits that share those confounds
overstate generalization; sensor- and site-dependent background structure is a
documented confound in bioacoustics \citep{lostanlen2019robust}, and random
cross-validation over spatially structured ecological data is known to
underestimate error \citep{roberts2017crossval}. Our diagnosis instantiates
this for forest threat detection: a leakage-free but random split hides a
cross-site collapse, and a linear probe shows the representation encodes site
identity. We use leave-one-site-out evaluation and false-positive-budgeted recall
to measure what random splits and accuracy conceal.

